\documentclass[11pt,a4paper]{article}

\usepackage[margin=2.3cm]{geometry}
\usepackage{array}
\usepackage{booktabs}
\usepackage{longtable}
\usepackage{hyperref}

\Urlmuskip=0mu plus 2mu
\setlength{\emergencystretch}{3em}
\newcommand{\field}[1]{\nolinkurl{#1}}
\newcommand{\code}[1]{\texttt{\detokenize{#1}}}

\title{CTR Diagnostic CSV File: Field Definitions and Computation}
\author{}
\date{}

\begin{document}
\sloppy
\maketitle

\section{Scope}

This note documents the CSV file written by \code{CTRDiagnosticLogger} in
\code{ctr_two_tubes_dynamic.py} and \code{ctr_two_tubes_quasi-static.py}.
The logger is attached to the two-tube concentric-tube-robot scene and writes
one row from \code{onAnimateEndEvent} every \code{every_n_steps} simulation
steps. In the current scenes this sampling period is 20 steps.

Tube 1, named \code{Tube_1} in the scene, is the outer tube. Tube 2 in the
diagnostic column names is the inner tube, named \code{Tube_3} in the scene.
The rigid pose vectors are SOFA \code{Rigid3d} states
\([x,y,z,q_x,q_y,q_z,q_w]\). The Cosserat strain states are SOFA
\code{Vec3d} arrays, one entry per section, and the third component
\code{state[i][2]} is the bending curvature \(\kappa\) used by the diagnostic
curvature fields.

\section{Common Computation Rules}

At construction time the logger stores the initial outer-tube base pose
\code{t1_base_mo.position.value[0]} and the initial outer-tube frame centers
from \code{t1_frames_mo.position.value}. All base and frame drift fields are
computed relative to those stored values.

For a Cosserat mechanical object \code{mo}, the helper
\code{_strain_stats(mo)} reads
\code{mo.position.value} as the current section strain vectors and
\code{mo.rest_position.value} as the rest strain vectors. For each section
\(i\), the strain error is
\[
  e_i = \left\| \mathbf{s}_i - \mathbf{s}^{rest}_i \right\|_2 .
\]
The reported strain errors are maxima or means of these section errors. The
reported curvature fields are means of the third strain component,
\(\kappa_i=s_i[2]\), over selected sections.

A ``curved'' section is any selected section whose rest curvature satisfies
\[
  \left|s^{rest}_i[2]\right| > 10^{-12}.
\]
The helper \code{_curved_kappa_mean} returns the mean current curvature and
the mean rest curvature over those curved sections. If no section is selected,
the logger writes zero.

Curvature relative-error columns are written as percentages:
\[
  100\,\frac{\left|\bar{\kappa}-\bar{\kappa}^{rest}\right|}
            {\left|\bar{\kappa}^{rest}\right|}.
\]
If the rest-curvature denominator is below \(10^{-12}\), the logger writes
zero to avoid division by a near-zero reference.

For the inner tube exposed-region fields, the logger first computes
\[
  \mathrm{t2\_advance} = x_{\mathrm{t2,base}} - x_{\mathrm{t2,offset}},
  \qquad
  \mathrm{protrusion} =
  \max\left(0,\mathrm{t2\_advance}-x_{\mathrm{t1,base}}\right).
\]
Then
\[
  s_{\mathrm{exposed,start}} =
  L_{\mathrm{t2}} - \mathrm{protrusion},
\]
where \(L_{\mathrm{t2}}\) is \code{T2_PARAMS["str_length"]}. A section is
considered exposed when the midpoint of its intrinsic arc-length interval,
computed from \code{t2_sec_curv_abs}, is greater than or equal to this start
coordinate.

Contact gaps are recomputed by the logger from
\code{BeamContactMapping.contactTriads} and the mapped contact point
mechanical object \code{contact_mo.position}. For contact index \(k\), the
two mapped points are read as \(q=\code{position[2*k]}\) and
\(p=\code{position[2*k+1]}\), and the normal \(n\) is read from the
corresponding contact triad. The signed gap is
\[
  g_k =
  \code{gapSign}\, (p-q)\cdot n .
\]
A pair is counted as ``active-like'' when \(g_k \leq\)
\code{activationTolerance}. This is a diagnostic contact-activity criterion,
not a separate SOFA constraint state.

The final contact-reaction columns differ slightly between the two scene
files:
\begin{itemize}
  \item In \code{ctr_two_tubes_dynamic.py}, the columns are named
  \field{active_impulse_pairs}, \field{normal_impulse_sum}, and
  \field{normal_impulse_max}. The logger reads
  \code{cpuc.normalContactImpulses.value}, takes absolute values, and
  estimates forces by dividing impulse sums and maxima by \code{dt}.
  \item In \code{ctr_two_tubes_quasi-static.py}, the corresponding columns are
  named \field{active_lambda_pairs}, \field{normal_lambda_sum}, and
  \field{normal_lambda_max}. The same SOFA data field is read, but in the
  first-order quasi-static solve the stored lambda is treated as already
  force-like, so the force columns are equal to the lambda sum and maximum.
\end{itemize}

\section{CSV Fields}

\small
\begin{longtable}{>{\raggedright\arraybackslash}p{0.26\textwidth}
                  >{\raggedright\arraybackslash}p{0.30\textwidth}
                  >{\raggedright\arraybackslash}p{0.34\textwidth}}
\toprule
\textbf{Field} & \textbf{Meaning} & \textbf{How it is read or computed} \\
\midrule
\endfirsthead
\toprule
\textbf{Field} & \textbf{Meaning} & \textbf{How it is read or computed} \\
\midrule
\endhead
\bottomrule
\endfoot

\field{step} &
Logger animation-step counter at the sampled row. &
\code{self.step} is incremented on every \code{onAnimateEndEvent}; a row is
written only when \code{step \% every_n_steps == 0}. \\

\field{time} &
Current SOFA simulation time in seconds. &
Read from \code{root_node.getTime()}. \\

\field{phase} &
GUI/controller phase, for example \code{initializing} or \code{control}. &
Read from \code{gui_bridge.snapshot()["phase"]}. \\

\field{dt} &
Current scene time step in seconds. &
Read from \code{root_node.dt.value} with a safe float conversion. \\

\field{translation_step_um} &
Per-simulation-step translation rate limit, in micrometres per step. &
Read from \code{gui_bridge.snapshot()["translation_step_m"]} and multiplied
by \(10^6\). \\

\field{t1_base_x} &
Outer tube simulated rigid-base \(x\) position in metres. &
Read as \code{t1_base_mo.position.value[0][0]}. \\

\field{t1_ctrl_x} &
Outer tube control/rest target \(x\) position in metres. &
Read as \code{t1_control_mo.position.value[0][0]}. This is the external
target used by the base \code{RestShapeSpringsForceField}. \\

\field{t1_lag_x} &
Outer tube axial lag between target and simulated base. &
Computed as \code{t1_ctrl_x - t1_base_x}. \\

\field{t1_base_trans_delta_m} &
Magnitude of outer base translational drift from its initial pose, in metres. &
Computed as the Euclidean norm of the difference between the current base
\([x,y,z]\) and the stored initial base \([x,y,z]\). \\

\field{t1_base_rot_delta_rad} &
Outer base angular drift from its initial orientation, in radians. &
Computed from the current and initial quaternions using
\(2\arccos(|q\cdot q_0|)\), with the dot product clamped to \([-1,1]\). \\

\field{t1_strain_err_max} &
Maximum outer-tube section strain error. &
Maximum of \(\|\mathbf{s}_i-\mathbf{s}^{rest}_i\|_2\) over all Tube 1
Cosserat sections read from \code{t1_coss_mo.position} and
\code{t1_coss_mo.rest_position}. \\

\field{t1_strain_err_mean} &
Mean outer-tube section strain error. &
Mean of the same Tube 1 section errors used for
\field{t1_strain_err_max}. \\

\field{t1_kappa_mean_curved_1pm} &
Mean current curvature of the curved part of the outer tube, in \(1/\mathrm{m}\). &
Mean of \code{t1_coss_mo.position.value[i][2]} over Tube 1 sections whose
rest curvature is nonzero. \\

\field{t1_kappa_rest_mean_curved_1pm} &
Mean rest curvature of the curved part of the outer tube, in \(1/\mathrm{m}\). &
Mean of \code{t1_coss_mo.rest_position.value[i][2]} over the same curved Tube
1 sections. \\

\field{t1_kappa_rel_error} &
Percentage difference between the current and rest mean curvature of the curved
outer-tube sections. &
Computed as
\code{100 * abs(t1_kappa - t1_rest_kappa) / abs(t1_rest_kappa)}, with zero
written when the rest mean curvature is too close to zero. \\

\field{t1_frame_delta_max_m} &
Maximum displacement of an outer-tube output frame from its initial center, in metres. &
For each frame in \code{t1_frames_mo.position.value}, the logger compares the
current \([x,y,z]\) center with the stored initial center and reports the
maximum norm. \\

\field{t1_frame_delta_mean_m} &
Mean displacement of outer-tube output frames from their initial centers, in metres. &
Mean of the same frame-center displacement norms used for
\field{t1_frame_delta_max_m}. \\

\field{t2_base_x} &
Inner tube simulated rigid-base \(x\) position in metres. &
Read as \code{t2_base_mo.position.value[0][0]}. \\

\field{t2_ctrl_x} &
Inner tube control/rest target \(x\) position in metres. &
Read as \code{t2_control_mo.position.value[0][0]}. This is the external
target used by the base \code{RestShapeSpringsForceField}. \\

\field{t2_lag_x} &
Inner tube axial lag between target and simulated base. &
Computed as \code{t2_ctrl_x - t2_base_x}. \\

\field{t2_advance_m} &
Inner tube insertion measured from its initial offset, in metres. &
Computed as \code{t2_base_x - t2_x_offset}. In the current geometry
\code{t2_x_offset} is \(-0.16\) m from \code{compute_concentric_offset}. \\

\field{estimated_protrusion_m} &
Estimated length of inner tube protruding beyond the outer tube tip, in metres. &
Computed as \code{max(0, t2_advance_m - t1_base_x)}. \\

\field{t2_strain_err_max_all} &
Maximum inner-tube section strain error over all sections. &
Maximum of \(\|\mathbf{s}_i-\mathbf{s}^{rest}_i\|_2\) over all Tube 2
Cosserat sections read from \code{t2_coss_mo.position} and
\code{t2_coss_mo.rest_position}. \\

\field{t2_strain_err_mean_all} &
Mean inner-tube section strain error over all sections. &
Mean of the same all-section Tube 2 strain errors used for
\field{t2_strain_err_max_all}. \\

\field{t2_strain_err_max_exposed} &
Maximum inner-tube strain error over exposed sections only. &
The logger first selects exposed section ids using
\code{T2_PARAMS["str_length"] - estimated_protrusion_m} and the section
midpoints from \code{t2_sec_curv_abs}; it then takes the maximum strain error
over those ids. \\

\field{t2_strain_err_mean_exposed} &
Mean inner-tube strain error over exposed sections only. &
Mean of the exposed-section Tube 2 strain errors used for
\field{t2_strain_err_max_exposed}. \\

\field{t2_kappa_mean_exposed_curved_1pm} &
Mean current curvature of exposed, naturally curved inner-tube sections, in
\(1/\mathrm{m}\). &
Mean of \code{t2_coss_mo.position.value[i][2]} over exposed section ids whose
rest curvature is nonzero. \\

\field{t2_kappa_rest_mean_exposed_curved_1pm} &
Mean rest curvature of exposed, naturally curved inner-tube sections, in
\(1/\mathrm{m}\). &
Mean of \code{t2_coss_mo.rest_position.value[i][2]} over the same exposed and
curved section ids. \\

\field{t2_kappa_rel_err_exposed_curved} &
Percentage difference between the current and rest mean curvature of exposed,
naturally curved inner-tube sections. &
Computed as
\code{100 * abs(exposed_kappa - exposed_rest_kappa) / abs(exposed_rest_kappa)},
with zero written when the rest mean curvature is too close to zero. In the
dynamic scene, the same exposed/protruded selection is stored in local variables
named \code{protruded_*}. \\

\field{t2_tip_x} &
Inner tube tip \(x\) coordinate in metres. &
Read from the last frame in \code{t2_frames_mo.position.value[-1][0]}; if no
frame exists, zero is written. \\

\field{t2_tip_y} &
Inner tube tip \(y\) coordinate in metres. &
Read from \code{t2_frames_mo.position.value[-1][1]}; if no frame exists, zero
is written. \\

\field{t2_tip_z} &
Inner tube tip \(z\) coordinate in metres. &
Read from \code{t2_frames_mo.position.value[-1][2]}; if no frame exists, zero
is written. \\

\field{contact_pairs} &
Number of valid contact gaps reconstructed by the logger. &
Counts contact triads from \code{bcm.contactTriads.value} that have a valid
normal and a matching pair of mapped points in \code{contact_mo.position}. \\

\field{active_like_pairs} &
Number of reconstructed contact pairs whose gap is within the activation
tolerance. &
Counts reconstructed gaps satisfying
\(g_k \leq \code{cpuc.activationTolerance.value}\). \\

\field{min_gap} &
Minimum reconstructed signed contact gap in metres. &
Minimum of \(g_k=\code{gapSign}(p-q)\cdot n\) over reconstructed gaps. If no
gap is reconstructed, zero is written. \\

\field{max_gap} &
Maximum reconstructed signed contact gap in metres. &
Maximum of the same reconstructed gap list. If no gap is reconstructed, zero
is written. \\

\field{mean_gap} &
Mean reconstructed signed contact gap in metres. &
Mean of the same reconstructed gap list. If no gap is reconstructed, zero is
written. \\

\field{gap_sign} &
Sign convention applied to contact gaps. &
Read from \code{bcm.gapSign.value}; default fallback is \(1\) if the read
fails. \\

\field{activation_tolerance} &
Gap threshold used by the unilateral contact constraint for activation-like
diagnostics. &
Read from \code{cpuc.activationTolerance.value}; in the current scene this is
\(10^{-4}\) m. \\

\field{active_impulse_pairs} / \field{active_lambda_pairs} &
Number of stored normal contact reaction entries. &
Length of \code{cpuc.normalContactImpulses.value}. The dynamic scene writes
the impulse-name column; the quasi-static scene writes the lambda-name column. \\

\field{normal_impulse_sum} / \field{normal_lambda_sum} &
Sum of absolute normal contact reactions. &
The logger converts \code{cpuc.normalContactImpulses.value} to floats, takes
absolute values, and sums them. Dynamic rows interpret this as impulse sum;
quasi-static rows interpret it as lambda or force-like reaction sum. \\

\field{normal_impulse_max} / \field{normal_lambda_max} &
Maximum absolute normal contact reaction. &
Maximum absolute value in \code{cpuc.normalContactImpulses.value}. Dynamic
rows interpret this as maximum impulse; quasi-static rows interpret it as
maximum lambda or force-like reaction. \\

\field{normal_force_sum_N} &
Estimated total normal contact force in newtons. &
In the dynamic scene, computed as
\field{normal_impulse_sum}/\field{dt} when \field{dt} is positive, otherwise
zero. In the quasi-static scene, equal to \field{normal_lambda_sum}. \\

\field{normal_force_max_N} &
Estimated maximum single-pair normal contact force in newtons. &
In the dynamic scene, computed as
\field{normal_impulse_max}/\field{dt} when \field{dt} is positive, otherwise
zero. In the quasi-static scene, equal to \field{normal_lambda_max}. \\

\end{longtable}
\normalsize

\section{Practical Debugging Interpretation}

Large \field{t1_lag_x} or \field{t2_lag_x} means the simulated base is not
following its spring-control target. Large \field{t1_base_trans_delta_m} or
\field{t1_base_rot_delta_rad} during initialization means the outer tube base
is moving even though it should be held near its initial pose.

For the inner tube, compare \field{t2_strain_err_mean_exposed} and
\field{t2_kappa_mean_exposed_curved_1pm} with
\field{t2_kappa_rest_mean_exposed_curved_1pm}. If the exposed section errors
stay large after protrusion, the inner tube is not relaxing toward its own
rest curvature. Contact columns then indicate whether this correlates with
near-active gaps or large normal reactions.

\end{document}
